\section{Symmetric Stable Products and Anisotropic Geometry}
\label{sec:symmetric-stable-products-and-anisotropic-geometry}

The finite-variance criterion does not extend unchanged to symmetric stable
products: both the coefficient domain and the mixing cost become
\(\alpha\)-dependent.
Fix \(0<\alpha<2\), let \(\rho_\alpha\) be the standard symmetric
\(\alpha\)-stable law, and put
\begin{equation}
 \int_\R e^{itx}\,\rho_\alpha(dx)=e^{-|t|^\alpha},
 \qquad
 \mu_\alpha=\rho_\alpha^{\otimes\N},
 \qquad
 \mu_{\alpha,n}=\rho_\alpha^{\otimes n}.
 \label{eq:standard-symmetric-stable-law}
\end{equation}
Let \((X_j)_{j\geq1}\) be independent with common law \(\rho_\alpha\), and
write \(\mathbf X=(X_j)_{j\geq1}\).
Its density \(f_\alpha\) is smooth and strictly positive, and its location
and scale Fisher information are finite, although its second moment is
infinite~\cite{bobkov2014stable}.
Indeed, \(f_\alpha(x)\asymp(1+|x|)^{-1-\alpha}\) and its location score is
\(O((1+|x|)^{-1})\) in the tails, so both the location score and its product
with \(x\) are square integrable against \(f_\alpha\).
For \(\alpha\le1\), the normalization in
\eqref{eq:standard-symmetric-stable-law} specifies a zero location parameter,
not a finite mean.

A stable row series \(\sum_ja_jX_j\) converges almost surely precisely when
\(a\in\ell^\alpha\).
For a real matrix \(T\) with \(\ell^\alpha\) rows, write
\(\Phi_{\alpha,T}\) for the associated row-series map and
\(\nu_{\alpha,T}=(\Phi_{\alpha,T})_\#\mu_\alpha\).
We say that \(T\) admits a compatible stable row-series realization if
there is a matrix \(S\)
with \(\ell^\alpha\) rows such that the two row-series maps are measurable
and
\[
 \Phi_{\alpha,S}\circ\Phi_{\alpha,T}=I,
 \qquad
 \Phi_{\alpha,T}\circ\Phi_{\alpha,S}=I
\]
on a common measurable carrier of full measure for the corresponding orbit
laws.
This is the stable analogue of the compatible realization supplied by
\cref{lem:measurable-row-series-realization}; it is an assumption on the
row-series realization, not a claim that every bounded \(\ell^2\) operator has
stable row sums.

For a matrix \(K\), define its column cost by
\[
 F_\alpha(K)=\sum_j\|Ke_j\|_2^\alpha.
\]
For \(M=(m_{ij})\), define the anisotropic mixing cost
\begin{equation}
 \mathfrak e_\alpha(M)
 =\sum_i|m_{ii}-1|^2
  +F_\alpha(M-\diag M).
 \label{eq:stable-mixing-cost}
\end{equation}
The first term measures coordinate scales.
The second takes the \(\ell^2\)-norm down each off-diagonal input
column and then its \(\alpha\)-th power across source coordinates.
For the actual stable output marginals, set
\begin{equation}
 E_{\alpha,n}(T)
 =-2\log\Aff\bigl(
       \mu_{\alpha,n},(\pi_n)_\#\nu_{\alpha,T}\bigr).
 \label{eq:stable-projected-hellinger-energy}
\end{equation}

\begin{theorem}[Complete criterion for symmetric stable coordinates]
\label{thm:complete-criterion-symmetric-stable-coordinates}
Let \(0<\alpha<2\) and suppose that \(T\) admits a compatible stable
row-series realization.  Then
\begin{equation}
 \boxed{\begin{aligned}
 \nu_{\alpha,T}\sim\mu_\alpha
 &\quad\Longleftrightarrow\quad
 T|_H\in\GL(H)\ \text{and}\
 \mathfrak e_\alpha(TQ^{-1})<\infty
       \ \text{for some }Q\in\mathcal P^{\pm}\\
 &\quad\Longleftrightarrow\quad
 \sup_nE_{\alpha,n}(T)<\infty,
 \end{aligned}}
 \label{eq:symmetric-stable-equivalence}
\end{equation}
Here \(\mathfrak e_\alpha\) is the anisotropic mixing cost from
\eqref{eq:stable-mixing-cost}, whose off-diagonal component is the column
cost \(F_\alpha\).
If these conditions fail, then
\(\nu_{\alpha,T}\perp\mu_\alpha\) and
\(E_{\alpha,n}(T)\uparrow\infty\).
Conversely, if \(T\in\GL(H)\) satisfies the summability condition in
\eqref{eq:symmetric-stable-equivalence}, that condition itself supplies the
compatible stable row-series realization for \(T\) and its inverse.
The exact linear stabilizer in this class is \(\mathcal P^{\pm}\).
\end{theorem}

\begin{proof}
There are two transfers from the finite-variance argument and two genuinely
stable steps.  The measure-class dichotomy and the diagonal scale analysis
transfer once the translation space is identified as \(\ell^2\).  The first
stable step is a dimension-free upper bound obtained from the Gaussian
scale-mixture representation.  Necessity is harder: a Fourier half-density
transfer recovers only Hilbert--Schmidt proximity, and a bounded test supported
on rare source events, followed by finite sign changes and coordinate cuts, is
needed to recover the sharper
\(\ell^\alpha(\ell^2)\) column cost.  We keep these roles separate below.

\emph{Dichotomy, energy, and the coefficient domain.}
Put \(g_\alpha=\sqrt{f_\alpha}\) and define
\[
 a_\alpha(u)=\int_\R g_\alpha(x)g_\alpha(x-u)\,dx.
\]
Finite positive location Fisher information, everywhere positivity of the
density, and decay of translations in \(L^2\) give
\begin{equation}
 1-a_\alpha(u)\asymp_\alpha\min\{u^2,1\}.
 \label{eq:stable-translation-affinity}
\end{equation}
Kakutani's theorem therefore identifies \(\ell^2\) as exactly the space of
translations under which \(\mu_\alpha\) is quasi-invariant.
If \(\nu_{\alpha,T}\sim\mu_\alpha\), translating before and after the
compatible stable row-series realization shows that the restrictions of \(T\)
and of the inverse matrix \(S\) in that realization map this translation space
into itself.
The closed graph theorem then gives \(T|_H\in\GL(H)\).

The compatible two-sided inverse makes every finite family of output rows
linearly independent.
It therefore has a nonsingular input minor, and the remaining stable inputs
contribute an independent convolution factor; hence every finite output
block has an everywhere-positive density.
The finite-marginal likelihood ratios consequently form the same mean-one
martingale as in
\cref{lem:finite-marginal-likelihood-ratios}.
The product measure is quasi-invariant and ergodic under finitely supported
translations, while its row-series image has the same properties under their
images by \(T\).
The equivalent--singular theorem for quasi-invariant ergodic measures allows
these two translation spaces to differ and therefore applies
\cite{okazaki1985dichotomy}.
It follows exactly as in \eqref{eq:marginal-affinity-limit} that
\begin{equation}
 \nu_{\alpha,T}\sim\mu_\alpha
 \quad\Longleftrightarrow\quad
 \sup_nE_{\alpha,n}(T)<\infty,
 \label{eq:stable-energy-dichotomy}
\end{equation}
with singularity and divergent energy on the other branch.

\emph{Sufficiency.}
The column cost satisfies
\begin{equation}
 F_\alpha(AK)\le\|A\|_{\mathrm{op}}^\alpha F_\alpha(K),
 \qquad
 F_\alpha(KD)\le\|D\|_{\mathrm{op}}^\alpha F_\alpha(K)
 \label{eq:stable-column-class-bounds}
\end{equation}
for bounded \(A\) and bounded diagonal \(D\), and
\begin{equation}
 \|K^*\gamma\|_\alpha^\alpha
 \le F_\alpha(K)\|\gamma\|_2^\alpha.
 \label{eq:stable-row-admissibility-bound}
\end{equation}
Thus \(F_\alpha(K)<\infty\) implies that \(K\) has admissible stable rows.

The finite-dimensional estimate that replaces the finite-variance Hellinger
path bound is
\begin{equation}
 h^2(A_\#\mu_{\alpha,n},\mu_{\alpha,n})
 \le C_\alpha\min\{\mathfrak e_\alpha(A),1\}
 \label{eq:stable-finite-dimensional-upper-bound}
\end{equation}
for every invertible real \(n\times n\) matrix \(A\), with a constant
independent of \(n\).
To prove it, put \(\beta=\alpha/2\), let \(V_i\) be independent positive
\(\beta\)-stable variables with
\(\E e^{-tV_i}=e^{-t^\beta}\), and write
\(X_i=\sqrt{2V_i}Z_i\) with independent standard Gaussians \(Z_i\).
Conditionally on the diagonal matrix \(V\) with entries \(V_i\), Hellinger
data processing and the
Gaussian affinity formula bound the left side of
\eqref{eq:stable-finite-dimensional-upper-bound} by a constant times
\[
 \E\min\{\|V^{-1/2}(A-I)V^{1/2}\|_F^2,1\}.
\]
For \(c_j^2=\sum_{i\ne j}|a_{ij}|^2\), the contribution of source column
\(j\) is
\[
 W_j=V_j\sum_{i\ne j}\frac{|a_{ij}|^2}{V_i}.
\]
The column estimate is the point at which the stable exponent enters.  The
positive stable estimates, for \(y>0\),
\[
 \mathbb P(V_1>y)\le C_\beta y^{-\beta},
 \qquad
 \E[V_1;V_1\le y]\le C_\beta y^{1-\beta},
 \qquad
 \E V_1^{-1}<\infty
\]
show that the contribution of \(\{c_j^2V_j>1\}\) is at most
\(C_\alpha c_j^{2\beta}=C_\alpha c_j^\alpha\).  For \(c_j>0\), independence
gives on the complementary event
\[
 \E[W_j;c_j^2V_j\le1]
 \le C_\beta c_j^2\E[V_j;V_j\le c_j^{-2}]
 \le C_\alpha c_j^\alpha.
\]
The zero-column case is immediate.  Thus source column \(j\) contributes at
most \(C_\alpha c_j^\alpha\) to this upper bound, regardless of how its mass is
distributed among receiver coordinates.  Summing these bounds, while retaining
the ordinary quadratic estimate for diagonal scales, proves
\eqref{eq:stable-finite-dimensional-upper-bound}.

Now remove the measure-preserving \(Q\) and write the resulting matrix as
\(M=D+K\), where \(D\) is diagonal with entries \(d_i\),
\((d_i-1)\in\ell^2\), and
\(F_\alpha(K)<\infty\).
The finitely many indices on which the diagonal of \(M\) is not positive
are absorbed into \(K\).
For
\[
 D_n=I+P_n(D-I)P_n,
 \qquad K_n=P_nKP_n,
 \qquad M_n=D_n+K_n,
\]
we have \(M_n\to M\) in operator norm, so the \(M_n\) are eventually
invertible with uniformly bounded inverses.
The identity
\[
 M_n^{-1}=D_n^{-1}-M_n^{-1}K_nD_n^{-1}
\]
and the column-class bounds show that
\(F_\alpha(M_n^{-1}K_nD_n^{-1})\) is uniformly bounded.  Moreover,
\begin{equation}
 \begin{aligned}
 M_n^{-1}M_m-I
 ={}&D_n^{-1}(D_m-D_n)-M_n^{-1}K_nD_n^{-1}(D_m-D_n)\\
    &+M_n^{-1}(K_m-K_n).
 \end{aligned}
 \label{eq:stable-relative-approximants}
\end{equation}
The first term has diagonal \(\ell^2\)-norm tending to zero.
The other two have column cost tending to zero by
\eqref{eq:stable-column-class-bounds},
\(\|D_m-D_n\|_{\mathrm{op}}\to0\), and
\(F_\alpha(K_m-K_n)\to0\); their diagonal square cost tends to zero as
well because finite column cost implies membership in \(\mathcal S_2\).
Thus \eqref{eq:stable-finite-dimensional-upper-bound}, applied to
\(M_n^{-1}M_m\), makes the square-root likelihoods of
\((M_n)_\#\mu_\alpha\) Cauchy in \(L^2(\mu_\alpha)\).
Their limit has norm one, and convergence of each stable row in
\(\ell^\alpha\) identifies its square with
\(d\nu_{\alpha,M}/d\mu_\alpha\).
This proves absolute continuity and then equivalence by
\eqref{eq:stable-energy-dichotomy}.
Finally,
\begin{equation}
 (D+K)^{-1}
 =D^{-1}-(D+K)^{-1}KD^{-1}
 \label{eq:stable-inverse-decomposition}
\end{equation}
and \eqref{eq:stable-column-class-bounds} show that the inverse has the same
admissible form.
The continuity of the stable random functional from \(\ell^\alpha\) to
convergence in probability lets the finite-matrix composition identities
pass to the row-series limits, producing compatible measurable versions.
This also proves the last assertion of the theorem.

\emph{Reduction to a Hilbert--Schmidt perturbation.}
Suppose conversely that \(\nu_{\alpha,T}\sim\mu_\alpha\).  We first extract
the quadratic information that equivalence still forces.  The auxiliary law
is dictated by the quantity already controlled by equivalence: the translation
affinity \(a_\alpha\) is the autocorrelation of the half-density \(g_\alpha\),
hence the characteristic function of the probability obtained by squaring the
modulus of its unitary Fourier transform.  Moreover, a linear change of
variables on half-densities becomes the adjoint inverse on the Fourier side.  A
convolution power will regularize this auxiliary law enough to apply the
finite-variance theorem without changing the intertwining relation.  This
transfer yields only a Hilbert--Schmidt conclusion; the stable column cost will
be recovered separately.

Let \(\widehat g_\alpha\) be the unitary Fourier transform of \(g_\alpha\),
and set
\[
 \lambda(d\xi)=|\widehat g_\alpha(\xi)|^2\,d\xi.
\]
The probability \(\lambda\) is symmetric and has variance
\(\|g_\alpha'\|_2^2\in(0,\infty)\).
Moreover, \(\widehat g_\alpha*\widehat g_\alpha\) is a positive constant multiple of
\(e^{-|t|^\alpha}\); hence
\(|\widehat g_\alpha|^2*|\widehat g_\alpha|^2>0\) almost everywhere.
The characteristic function of \(\lambda\) is \(a_\alpha\).
Polynomial stable tails give, for some \(m_0,\delta_0>0\),
\[
 c(1+|t|)^{-m_0}\le a_\alpha(t)
 \le C(1+|t|)^{-\delta_0},
\]
and \(a_\alpha'\) is bounded.
Choose an integer \(r\) so large that
\[
 \int_\R t^2\bigl(|a_\alpha(t)|^{2r}
                 +|(a_\alpha^r)'(t)|^2\bigr)\,dt<\infty.
\]
Fourier inversion makes the density of \(\lambda^{*r}\) a bounded-variation
density; three such convolutions have finite location Fisher information by
the Fisher smoothing estimate~\cite{bobkov2014stable}.
Let \(f_\lambda\) be the density of \(\lambda^{*3r}\), put
\(f_{\lambda,2}=f_\lambda*f_\lambda\), and let
\(s_\lambda,s_{\lambda,2}\) be their location scores.
Independent \(U_1,U_2\) with density \(f_\lambda\) satisfy
\[
 x s_{\lambda,2}(x)-1
 =2\E[U_1s_\lambda(U_2)\mid U_1+U_2=x].
\]
Their finite variance and the finite Fisher information of \(f_\lambda\)
therefore
give finite scale Fisher information for \(f_{\lambda,2}\).
Set \(k=6r\) and \(\Lambda_k=(\lambda^{*k})^{\otimes\N}\).
The density of \(\lambda^{*k}\), after
rescaling to variance one, is a.e. positive and satisfies
\eqref{eq:main-assumptions}.
It is non-Gaussian because its characteristic function
\(a_\alpha^k\) has a polynomial lower bound, whereas a nondegenerate
Gaussian characteristic function has Gaussian decay.

Let \(f_{\alpha,n}=f_\alpha^{\otimes n}\), let \(f_{\alpha,T,n}\) be the
density of the true first \(n\) stable outputs, and put
\[
 \ell_{\alpha,T,n}=\frac{f_{\alpha,T,n}}{f_{\alpha,n}},
 \qquad
 \ell_{\alpha,T}=\frac{d\nu_{\alpha,T}}{d\mu_\alpha}.
\]
Here \(\ell_{\alpha,T,n}\) is lifted to the product space and
\(\sqrt{f_{\alpha,T,n}}
  =g_\alpha^{\otimes n}\sqrt{\ell_{\alpha,T,n}}\).
Equivalence gives
\(\sqrt{\ell_{\alpha,T,n}}\to\sqrt{\ell_{\alpha,T}}\)
in \(L^2(\mu_\alpha)\).

The construction is easiest to read backwards.  We seek a probability whose
characteristic function at finitely supported \(z\) is
\(\prod_i a_\alpha((T^{-1}z)_i)^k\), because that is the characteristic
function of \(((T^{-1})^*)_\#\Lambda_k\).  Translation of a half-density
becomes phase multiplication after Fourier transform, so one copy produces
one autocorrelation factor; \(k\) tensor copies produce the \(k\)-th power;
and addition of the \(k\) frequencies at each site realizes convolution by
\(\lambda^{*k}\).  Thus each operation below corresponds to one factor in the
desired intertwining identity.

For each \(n\), take \(k\) tensor copies of
\(\sqrt{f_{\alpha,T,n}}\), apply the
finite-dimensional unitary Fourier transform, square the modulus, and push
the resulting probability on \(\R^{kn}\) forward by adding the \(k\)
frequency coordinates belonging to each site.
Call the resulting probability on \(\R^n\) \(\xi_{\alpha,T,n}\), and extend it to
the full sequence space by
\begin{equation}
 \Xi_{\alpha,T,n}=\xi_{\alpha,T,n}\otimes
       (\lambda^{*k})^{\otimes\{i>n\}}.
 \label{eq:stable-fourier-transfer-approximants}
\end{equation}
The finite-dimensional pushforward \(\xi_{\alpha,T,n}\) has a Lebesgue density, and
the density of \(\lambda^{*k}\) is positive almost everywhere; hence
\(\Xi_{\alpha,T,n}\ll\Lambda_k\).

For \(m>n\), pad \(\sqrt{f_{\alpha,T,n}}\) by \(m-n\) factors of
\(g_\alpha\) before making the Fourier comparison.
Using total variation in the \(L^1\)-norm convention, unitarity, the inequality
\(\||u|^2-|v|^2\|_1\le2\|u-v\|_2\) for unit vectors, and contraction of
total variation under pushforward give
\begin{equation}
 \|\Xi_{\alpha,T,m}-\Xi_{\alpha,T,n}\|_{\mathrm{TV}}
 \le2\|(\sqrt{\ell_{\alpha,T,m}})^{\otimes k}
           -(\sqrt{\ell_{\alpha,T,n}})^{\otimes k}\|_{L^2(\mu_\alpha^{\otimes k})}
 \le2k\|\sqrt{\ell_{\alpha,T,m}}
           -\sqrt{\ell_{\alpha,T,n}}\|_{L^2(\mu_\alpha)}.
 \label{eq:stable-fourier-transfer-tv-bound}
\end{equation}
Thus \(\Xi_{\alpha,T,n}\) converges in total variation to a probability
\(\Xi_{\alpha,T}\ll\Lambda_k\).

Let \(\nu_{\alpha,T}^{[n]}=\ell_{\alpha,T,n}\mu_\alpha\); this measure has the true first \(n\)
output marginal and the original independent tail.
For every finitely supported \(z\), Fourier autocorrelation and
the square-root likelihood convergence give
\begin{align*}
 \widehat\Xi_{\alpha,T}(z)
 &=\lim_n\Aff\bigl(\nu_{\alpha,T}^{[n]},
             (x\mapsto x+z)_\#\nu_{\alpha,T}^{[n]}\bigr)^k\\
 &=\Aff\bigl(\nu_{\alpha,T},
        (x\mapsto x+z)_\#\nu_{\alpha,T}\bigr)^k\\
 &=\prod_i a_\alpha((T^{-1}z)_i)^k.
\end{align*}
The last expression is the characteristic function of
\(((T^{-1})^*)_\#\Lambda_k\).  Agreement for every finitely supported \(z\)
identifies every finite marginal, and hence the probability on the sequence
space:
\(\Xi_{\alpha,T}=((T^{-1})^*)_\#\Lambda_k\).
Consequently
\begin{equation}
 ((T^{-1})^*)_\#\Lambda_k\ll\Lambda_k.
 \label{eq:stable-fourier-transfer}
\end{equation}
The equivalent--singular alternative in \cref{thm:measure-equivalence}
upgrades this absolute continuity to equivalence.  Applying that theorem to
the product of the \(\lambda^{*k}\) law rescaled to variance one then gives
\(Q\in\mathcal P^{\pm}\) such that
\((T^{-1})^*-Q\in\mathcal S_2\).
Since
\[
 T-Q=-T\bigl((T^{-1})^*-Q\bigr)^*Q,
\]
we obtain \(T-Q\in\mathcal S_2\).
After removing \(Q\), write
\begin{equation}
 TQ^{-1}=I+K,
 \qquad K\in\mathcal S_2.
 \label{eq:stable-hilbert-schmidt-match}
\end{equation}

\emph{Necessity of the off-diagonal column summability.}
The Fourier transfer sees only \(K\in\mathcal S_2\).  For \(\alpha<2\) this
is too weak: when \(0<c_j=\|Be_j\|_2<1\), the required cost
\(c_j^\alpha\) is larger than \(c_j^2\).  The missing mass is detected on the
rare event \(|X_j|\asymp c_j^{-1}\): it has probability of order
\(c_j^\alpha\), while the receiver displacement \(X_jBe_j\) then has
order-one norm.  The test below gates on precisely these events.  A tail
estimate controls interference from the other columns, a
directional statistic prevents cancellation, and the probability of two open
gates makes the truncation error quadratic in the total cost.  These three
ingredients prove the following dimension-free estimate, valid for every
dimension and every real square matrix \(B\): there are
\(\varepsilon_\alpha,c_\alpha>0\) such that
\begin{equation}
 \diag B=0,\quad F_\alpha(B)\le\varepsilon_\alpha
 \quad\Longrightarrow\quad
 h^2((I+B)_\#\mu_{\alpha,n},\mu_{\alpha,n})
 \ge c_\alpha F_\alpha(B).
 \label{eq:stable-off-diagonal-column-lower-bound}
\end{equation}
To prove it, write \(F=F_\alpha(B)\),
\(c_j=\|Be_j\|_2\), and
\(u_j=Be_j/c_j\) when \(c_j\ne0\).
Then \(\|u_j\|_2=1\), \((u_j)_j=0\), and the stable scale contributed to
output row \(i\) satisfies
\begin{equation}
 \omega_i(B):=\sum_j|b_{ij}|^\alpha\le F.
 \label{eq:stable-row-scale-bound}
\end{equation}

First, for any finite family of vectors \(v_j\) in a real Hilbert space,
\begin{equation}
 \mathbb P\left(\left\|\sum_jv_jX_j\right\|_2>t\right)
 \le C_\alpha t^{-\alpha}\sum_j\|v_j\|_2^\alpha.
 \label{eq:stable-hilbert-random-sum-tail}
\end{equation}
Indeed, exclude the events
\(\|v_j\|_2|X_j|>t\), truncate on their complement, and apply Chebyshev's
inequality to the centered truncated sum.
The stable tail and truncated second moment give
\eqref{eq:stable-hilbert-random-sum-tail}.
In particular, there is a deterministic modulus
\(\tau_\alpha(s)\downarrow0\) as \(s\downarrow0\) such that
\begin{equation}
 \E\min\left\{\left\|\sum_jv_jX_j\right\|_2,1\right\}
 \le\tau_\alpha\left(\sum_j\|v_j\|_2^\alpha\right).
 \label{eq:stable-hilbert-random-sum-truncation}
\end{equation}

Second, one bounded statistic detects a unit displacement in every
Euclidean direction, uniformly in the dimension.
Take \(\psi(x)=\tanh x\).
There are a smooth odd increasing function \(\chi\), bounded in modulus by
one, and
\(c_{\alpha,\mathrm{rec}}>0\), depending only on \(\rho_\alpha\), such that for every unit
vector \(u\),
\begin{equation}
 \Theta_u(x)=\sum_i u_i\psi(x_i),
 \qquad
 \mathcal M_u(t)=\E\chi(\Theta_u(\mathbf X+tu)),
 \qquad
 \mathcal M_u(t)\ge c_{\alpha,\mathrm{rec}}\quad(1\le t\le2).
 \label{eq:stable-directional-detection}
\end{equation}
Here uniformity in \(u\) requires care.  Since \(\psi\) is odd and bounded,
\(\E\Theta_u(\mathbf X)=0\) and
\(\operatorname{Var}\Theta_u(\mathbf X)
=\operatorname{Var}(\psi(X))\sum_i u_i^2\), so one compact interval contains
all but a prescribed small amount of its mass.  Also
\[
 |\Theta_u(\mathbf X+tu)-\Theta_u(\mathbf X)|
 \le |t|\sum_i u_i^2\le2.
\]
Choose \(\chi'\) bounded below by a positive constant on that interval enlarged
by two on each side.
Since
\[
 \inf_{|\zeta|\le2}\E\psi'(X+\zeta)>0,
\]
the expectation of
\(\sum_i u_i^2\psi'(X_i+tu_i)\) is bounded below by this positive infimum.
Choose the interval so that the discarded contribution is smaller than half
that infimum.  Differentiation then gives a dimension-free positive lower
bound for
\[
 \mathcal M_u'(t)=\E\left[
   \chi'(\Theta_u(\mathbf X+tu))\sum_i u_i^2\psi'(X_i+tu_i)
                 \right]
\]
when \(|t|\le2\).  Oddness then proves
\eqref{eq:stable-directional-detection}.

Third, choose a nonzero smooth odd function \(w\), supported where
\(1\le|t|\le2\), with \(0\le w\le1\) on the positive half-line.
Put
\[
 \Theta_j(x)=\sum_{i\ne j}(u_j)_i\psi(x_i),
\]
and define the bounded test
\begin{equation}
 \Psi_B(x)=\max\left\{-1,\min\left\{
 \sum_{j:c_j\ne0}w(c_jx_j)\chi(\Theta_j(x)),1\right\}\right\}.
 \label{eq:stable-column-detection-test}
\end{equation}
For all sufficiently small \(c_j\), stable tail asymptotics give
\begin{equation}
 \mathbb P(1\le|c_jX_j|\le2)\le C_\alpha c_j^\alpha,
 \qquad
 \E|w(c_jX_j)|\ge c_\alpha c_j^\alpha.
 \label{eq:stable-gate-probability}
\end{equation}
The gates depend on independent input coordinates.
If \(N_{\mathrm{open}}\) is the number of open gates, then
\(\E[N_{\mathrm{open}}(N_{\mathrm{open}}-1)]\le C_\alpha F^2\).
Before the outer truncation, each summand in
\eqref{eq:stable-column-detection-test} has mean zero because its gate variable
is independent of its receiver variables.
The truncation changes the sum only when at least two gates are open, so
\begin{equation}
 |\E\Psi_B(\mathbf X)|\le C_\alpha F^2.
 \label{eq:stable-column-detection-input-mean}
\end{equation}

It remains to control the same statistic at
\(\mathbf Y=(I+B)\mathbf X\).
For fixed \(0<\delta<1/4\),
\eqref{eq:stable-row-scale-bound} and the stable tail bound imply
\begin{equation}
 \mathbb P\{\exists i:\ |c_i(B\mathbf X)_i|>\delta\}
 \le C_\alpha\delta^{-\alpha}\sum_i c_i^\alpha\omega_i(B)
 \le C_\alpha\delta^{-\alpha}F^2.
 \label{eq:stable-output-gate-exception}
\end{equation}
On the complementary event, Lipschitz continuity of \(w\) and
\eqref{eq:stable-gate-probability} show that replacing every output
gate \(w(c_jY_j)\) by its input gate \(w(c_jX_j)\) changes the expectation
by at most
\begin{equation}
 C_\alpha\delta F+C_\alpha\delta^{-\alpha}F^2.
 \label{eq:stable-output-gate-error}
\end{equation}
After this replacement, removing the outer truncation costs at most another
\(C_\alpha F^2\).

Finally, the random vector
\(B(I-e_j\otimes e_j)\mathbf X\) is independent of the gate variable \(X_j\),
and \eqref{eq:stable-hilbert-random-sum-truncation} shows that replacing
\(\Theta_j(\mathbf Y)\) by
\(\Theta_j(\mathbf X+c_jX_ju_j)\) costs, after summing over \(j\), at most
\begin{equation}
 C_\alpha F\tau_\alpha(F).
 \label{eq:stable-receiver-interference-error}
\end{equation}
Because \((u_j)_j=0\), conditioning on \(X_j\) and using
\cref{eq:stable-directional-detection,eq:stable-gate-probability} gives
\begin{equation}
 \sum_j\E\bigl[
   w(c_jX_j)\chi(\Theta_j(\mathbf X+c_jX_ju_j))\bigr]
 \ge c_\alpha F.
 \label{eq:stable-column-detection-main-response}
\end{equation}
Choose \(\delta\) first so that the \(\delta F\) error is at most one quarter
of the main response.  With this \(\delta\) fixed, choose \(F\) so small that
the \(F^2\), \(\delta^{-\alpha}F^2\), and
\(F\tau_\alpha(F)\) errors together cost at most another quarter.
Equations
\cref{eq:stable-column-detection-input-mean,eq:stable-output-gate-error,eq:stable-receiver-interference-error,eq:stable-column-detection-main-response}
give
\[
 \E\Psi_B(\mathbf Y)-\E\Psi_B(\mathbf X)\ge c_\alpha F.
\]
The test can be nonzero only when a gate is open.
The coordinates of \(\mathbf Y\) have stable scales bounded by
\((1+F)^{1/\alpha}\), so
\[
 \E\Psi_B(\mathbf X)^2+\E\Psi_B(\mathbf Y)^2\le C_\alpha F.
\]
The Hellinger Cauchy--Schwarz inequality now gives
\(h^2((I+B)_\#\mu_{\alpha,n},\mu_{\alpha,n})\ge c_\alpha F\),
which is \eqref{eq:stable-off-diagonal-column-lower-bound}.

We use this estimate together with uniform invariance under finite tail sign
changes, thereby avoiding any triangularity assumption.
The local estimate alone cannot rule out \(F_\alpha(B)=\infty\), because it
applies only at small cost.  A finite sign change isolates the entries crossing
a coordinate cut.  Equivalence makes every sufficiently remote finite sign
change uniformly invisible in Hellinger distance, whereas the local estimate
makes any cut of a fixed small cost uniformly visible.  Random cuts supply
arbitrarily large crossing cost, and the vanishing cost of a single remote
coordinate lets us add coordinates until a prescribed small threshold is
first crossed.  The contradiction between invisibility and visibility is the
global step.

Let \(B\) have \(\ell^\alpha\) rows, zero diagonal,
\(\|B\|_{\mathrm{op}}\le1/10\), and belong to \(\mathcal S_2\).  Suppose
\(I+B\) admits a compatible stable row-series realization and
\((I+B)_\#\mu_\alpha\sim\mu_\alpha\).
For a finite coordinate set \(\mathcal J\), let \(Q_{\mathcal J}\) change
the signs in \(\mathcal J\)
and set
\[
 \Delta_{\mathcal J}=Q_{\mathcal J}BQ_{\mathcal J}-B,
 \qquad
 A_{\mathcal J}=(I+B)^{-1}Q_{\mathcal J}(I+B)Q_{\mathcal J}.
\]
The matrix \(\Delta_{\mathcal J}\) contains precisely the entries crossing
the cut between \(\mathcal J\) and its complement.
Because \(\mathcal J\) is finite, these entries form a sum of finitely many
\(\ell^2\) columns and finitely many \(\ell^\alpha\) rows; hence
\(F_\alpha(\Delta_{\mathcal J})<\infty\).
The finite-dimensional bound
\eqref{eq:stable-off-diagonal-column-lower-bound} extends to such infinite matrices by
applying it to \(P_n\Delta_{\mathcal J}P_n\) and using the Hellinger convergence from
the already proved sufficiency argument.

Put
\[
 K_{\mathcal J}=(I+B)^{-1}\Delta_{\mathcal J},
 \qquad
 D_{\mathcal J}=I+\diag K_{\mathcal J},
 \qquad
 \widetilde B_{\mathcal J}=D_{\mathcal J}^{-1}(I+K_{\mathcal J})-I.
\]
Then \(A_{\mathcal J}=I+K_{\mathcal J}\) and
\(\diag\widetilde B_{\mathcal J}=0\).
Since \(\|(I+B)^{-1}-I\|_{\mathrm{op}}\le1/9\), orthogonal projection
away from the diagonal entry in each column gives
\[
 \frac89\|\Delta_{\mathcal J}e_j\|_2
 \le\|(I-e_j\otimes e_j)K_{\mathcal J}e_j\|_2
 \le\frac{10}{9}\|\Delta_{\mathcal J}e_j\|_2.
\]
For small \(F_\alpha(\Delta_{\mathcal J})\), the diagonal entries of
\(D_{\mathcal J}\) lie in
\([1/2,3/2]\), and consequently
\begin{equation}
 F_\alpha(\widetilde B_{\mathcal J})\asymp_\alpha F_\alpha(\Delta_{\mathcal J}),
 \qquad
 \|\diag K_{\mathcal J}\|_2
 \le\frac19F_\alpha(\Delta_{\mathcal J})^{1/\alpha}.
 \label{eq:stable-finite-cut-comparison}
\end{equation}
The lower bound for \(\widetilde B_{\mathcal J}\), the diagonal scale upper
bound for \(D_{\mathcal J}\), and the Hellinger triangle inequality now give
\begin{equation}
 F_\alpha(\Delta_{\mathcal J})\le\varepsilon_*
 \quad\Longrightarrow\quad
 h((A_{\mathcal J})_\#\mu_\alpha,\mu_\alpha)
 \ge c_*F_\alpha(\Delta_{\mathcal J})^{1/2}.
 \label{eq:stable-finite-cut-lower-bound}
\end{equation}
Indeed, by \eqref{eq:stable-off-diagonal-column-lower-bound}, the two terms
have respective orders
\(F_\alpha(\Delta_{\mathcal J})^{1/2}\) and
\(F_\alpha(\Delta_{\mathcal J})^{1/\alpha}\), and the latter is smaller for small
cost because \(\alpha<2\).

On the other hand, finite-prefix approximation of the square-root likelihood
gives uniform invariance under finite tail sign changes.  Every tail sign
change fixes a prefix-dependent approximation and preserves \(\mu_\alpha\);
the relevant Hellinger distance is at most twice the \(L^2\)-approximation
error, uniformly over all finite tail sign sets.  Hence
\begin{equation}
 \sup_{\substack{\mathcal J\subset\{N+1,N+2,\ldots\}\\ |\mathcal J|<\infty}}
 h\bigl((I+B)_\#\mu_\alpha,
        (Q_{\mathcal J})_\#(I+B)_\#\mu_\alpha\bigr)
 \longrightarrow0.
 \label{eq:stable-tail-sign-change-invariance}
\end{equation}
Hellinger invariance and \((Q_{\mathcal J})_\#\mu_\alpha=\mu_\alpha\) also give the
connecting identity
\begin{equation}
 h\bigl((I+B)_\#\mu_\alpha,
        (Q_{\mathcal J})_\#(I+B)_\#\mu_\alpha\bigr)
 =h(\mu_\alpha,(A_{\mathcal J})_\#\mu_\alpha).
 \label{eq:stable-finite-cut-distance-identity}
\end{equation}
If \(F_\alpha(B)=\infty\), deleting finitely many columns changes the cost
by a finite amount, and so does deleting finitely many rows because those
rows lie in \(\ell^\alpha\).
Hence every tail still has infinite column cost.
Let \(I_0\) be a finite coordinate set in such a tail, and
form \(\mathcal J\subset I_0\) by independently including each
coordinate with probability
one half.
Column by column, concavity of \(t^{\alpha/2}\) gives
\begin{equation}
 \E_{\mathcal J}F_\alpha(P_{I_0}\Delta_{\mathcal J}P_{I_0})
 \ge2^{\alpha-1}F_\alpha(P_{I_0}BP_{I_0}).
 \label{eq:stable-random-finite-cut}
\end{equation}
As \(I_0\) increases through finite subsets of the tail, the right side is
unbounded, so some finite cuts have arbitrarily large cost.
Equivalence also makes every sufficiently remote one-coordinate output
marginal close to \(\rho_\alpha\), so
\[
 \omega_i(B):=\sum_{j\ne i}|b_{ij}|^\alpha\longrightarrow0.
\]
Together with \(B\in\mathcal S_2\), which gives
\(c_i:=\|Be_i\|_2\to0\), this yields
\[
 F_\alpha(\Delta_{\{i\}})
 =2^\alpha(c_i^\alpha+\omega_i(B))\longrightarrow0.
\]
For \(\alpha\le1\), \(F_\alpha\) is subadditive; for \(\alpha\ge1\),
\(F_\alpha^{1/\alpha}\) is the \(\ell^\alpha(\ell^2)\) norm and satisfies
the triangle inequality.
Moreover, when \(i\notin\mathcal J\),
\[
 \Delta_{\mathcal J\cup\{i\}}-\Delta_{\mathcal J}
 =Q_{\mathcal J}\Delta_{\{i\}}Q_{\mathcal J},
\]
and sign multiplication preserves the relevant subadditive size.
Thus the increment from adding one sufficiently remote coordinate tends
uniformly to zero in that size.
Fix \(0<\varepsilon<\varepsilon_*/2\).
Starting with the empty cut and adding the elements of a finite cut supplied
by \eqref{eq:stable-random-finite-cut}, stop the first time the corresponding
threshold is crossed.
For every sufficiently remote tail this produces a finite \(\mathcal J\) with
\(\varepsilon\le F_\alpha(\Delta_{\mathcal J})\le2\varepsilon\).
Equations
\cref{eq:stable-finite-cut-lower-bound,eq:stable-tail-sign-change-invariance,eq:stable-finite-cut-distance-identity} then contradict one another.
Hence \(F_\alpha(B)<\infty\).

Finally, put \(M=I+K\) in
\eqref{eq:stable-hilbert-schmidt-match} and let
\(K_n=P_nKP_n\) and \(M_n=I+K_n\).
For a sufficiently large \(n\), the equivalent transformation
\(A=M_n^{-1}M\) is arbitrarily close to the identity in operator norm.
Indeed, \(M_n\) changes only finitely many coordinates, so
\((M_n)_\#\mu_\alpha\sim\mu_\alpha\); compatible composition with
\(M_\#\mu_\alpha\sim\mu_\alpha\) gives
\(A_\#\mu_\alpha\sim\mu_\alpha\).
Its diagonal part \(D=\diag A\) is positive, boundedly invertible, and
\(D-I\in\mathcal S_2\); the one-coordinate scale criterion therefore gives
\(D_\#\mu_\alpha\sim\mu_\alpha\).
Consequently
\[
 D^{-1}A=I+B,
 \qquad
 \diag B=0,\quad B\in\mathcal S_2,\quad
 \|B\|_{\mathrm{op}}\le1/10,\quad
 (I+B)_\#\mu_\alpha\sim\mu_\alpha.
\]
Finite-coordinate and bounded diagonal composition preserve the
\(\ell^\alpha\) row condition and the compatible stable row-series
realization.  Thus the normalized \(B\) satisfies the hypotheses of the
preceding finite-cut contradiction, which gives \(F_\alpha(B)<\infty\).
Since \(K_n\) is supported on a finite coordinate block, we have
\[
 M=M_nD(I+B)
  =D+\bigl(K_nD+DB+K_nDB\bigr).
\]
The term in parentheses has finite column cost by
\eqref{eq:stable-column-class-bounds} and finite-coordinate support.
Restoring the initial signed permutation therefore gives
\[
 \sum_j\left(\sum_{i\ne j}|k_{ij}|^2\right)^{\alpha/2}<\infty,
\]
while \(K\in\mathcal S_2\) already gives
\(\sum_i|k_{ii}|^2<\infty\).
This is precisely \(\mathfrak e_\alpha(TQ^{-1})<\infty\), and completes
the necessity.

If \(\nu_{\alpha,T}=\mu_\alpha\), comparison of the joint characteristic
functions gives
\(\|T^*u\|_\alpha=\|u\|_\alpha\) for every finitely supported \(u\).
The same statement for the compatible inverse makes \(T^*\) a surjective
isometry of \(\ell^\alpha\).
For \(0<\alpha<2\), the equality case of
\[
 2\|u\|_\alpha^\alpha+2\|v\|_\alpha^\alpha
 -\|u+v\|_\alpha^\alpha-\|u-v\|_\alpha^\alpha\ge0
\]
shows that disjointly supported vectors have disjointly supported images.
The images of the coordinate vectors therefore have pairwise disjoint
supports; surjectivity makes each support a singleton, and norm preservation
makes its nonzero entry \(\pm1\).
Thus \(T\in\mathcal P^{\pm}\), while every member of \(\mathcal P^{\pm}\)
preserves \(\mu_\alpha\) by coordinate independence and symmetry.  This proves
the assertion about the exact linear
stabilizer and completes the proof.
\end{proof}

\subsection{Anisotropic stable Hellinger coercivity and sharpness}
\label{sec:anisotropic-stable-hellinger-coercivity-and-sharpness}

The two summands in \eqref{eq:stable-mixing-cost} have genuinely different
homogeneities.  Let \(\mathcal P_n^{\pm}\) denote the group of all
\(n\times n\) signed permutation matrices.  For a finite matrix \(A\), put
\[
 \mathfrak e_{\alpha,n}(A)
 =\min_{Q\in\mathcal P_n^{\pm}}\mathfrak e_\alpha(AQ^{-1}).
\]

\begin{proposition}[Anisotropic stable Hellinger coercivity]
\label{prop:anisotropic-stable-hellinger-coercivity}
Let \(A\in\GL(n,\R)\).  If all singular values of \(A\) lie in a fixed interval
\([a,b]\subset(0,\infty)\), then constants depending only on
\(\alpha,a,b\), and not on the dimension, satisfy
\begin{equation}
 c_{\alpha,a,b}\min\{\mathfrak e_{\alpha,n}(A),1\}
 \le h^2(A_\#\mu_{\alpha,n},\mu_{\alpha,n})
 \le C_{\alpha,a,b}\min\{\mathfrak e_{\alpha,n}(A),1\}.
 \label{eq:stable-hellinger-coercivity}
\end{equation}
\end{proposition}

\begin{proof}
The upper bound is \eqref{eq:stable-finite-dimensional-upper-bound}, after
removing the minimizing signed permutation.
For the even convolution power \(k\) constructed in the proof of
\cref{thm:complete-criterion-symmetric-stable-coordinates}, the
finite-dimensional Fourier transform gives the quantitative inequality
\begin{equation}
 h^2\bigl(((A^{-1})^*)_\#(\lambda^{*k})^{\otimes n},
                 (\lambda^{*k})^{\otimes n}\bigr)
 \le k\,h^2(A_\#\mu_{\alpha,n},\mu_{\alpha,n}).
 \label{eq:stable-finite-fourier-transfer}
\end{equation}
Indeed, the unitary Fourier transform intertwines the half-density actions
of \(A\) and \((A^{-1})^*\); taking moduli is an \(L^2\) contraction.
For \(k\) tensor copies, squared Hellinger distance grows by at most the
factor \(k\), and coordinatewise frequency addition is a data-processing
map.
Apply the finite-variance coercivity theorem to the product of the
\(\lambda^{*k}\) law rescaled to variance one.
When the singular values of \(A\) lie in \([a,b]\), those of \((A^{-1})^*\)
lie in \([b^{-1},a^{-1}]\), and
\[
 A-Q=-A\bigl((A^{-1})^*-Q\bigr)^*Q
\]
shows that distance from the signed-permutation group changes by at most a
factor depending on \(a,b\).
Thus \eqref{eq:stable-finite-fourier-transfer} gives
\begin{equation}
 h^2(A_\#\mu_{\alpha,n},\mu_{\alpha,n})
 \ge c_{\alpha,a,b}
     \min\left\{\min_{Q\in\mathcal P_n^{\pm}}\|A-Q\|_F^2,1\right\}.
 \label{eq:stable-quadratic-lower-bound}
\end{equation}

Near a signed permutation, remove that permutation and write
\[
 d_{\mathrm{diag}}=\sum_i|a_{ii}-1|^2,
 \qquad
 F=\sum_j\left(\sum_{i\ne j}|a_{ij}|^2\right)^{\alpha/2}.
\]
For \(d_{\mathrm{diag}}+F\) small, the identity is the nearest signed
permutation, so \eqref{eq:stable-quadratic-lower-bound} controls
\(d_{\mathrm{diag}}\).
If \(D=\diag A\) and \(D^{-1}A=I+B\), then
\(\diag B=0\) and \(F_\alpha(B)\asymp_\alpha F\).
The off-diagonal column lower bound and the diagonal scale upper bound give
\[
 h(A_\#\mu_{\alpha,n},\mu_{\alpha,n})
 \ge c_\alpha\sqrt F-C_\alpha\sqrt{d_{\mathrm{diag}}}.
\]
If \(d_{\mathrm{diag}}\) dominates \(F\), use
\eqref{eq:stable-quadratic-lower-bound}; otherwise absorb the second term in
the last display.
This proves the local lower bound in
\eqref{eq:stable-hellinger-coercivity}.
If uniform separation failed away from that neighborhood, one could take a
block-diagonal sum of counterexamples whose squared Hellinger distances are
summable.
Kakutani's theorem would make the resulting infinite stable products
equivalent.
The Hilbert--Schmidt part of
\cref{thm:complete-criterion-symmetric-stable-coordinates} forces the signed
permutation furnished by that theorem to preserve every sufficiently late block;
otherwise each crossed block contributes a fixed positive square error.
The restrictions of that permutation to the late blocks would then make
their anisotropic mixing costs summable, a contradiction.
This proves the global lower bound.
\end{proof}

The different local powers are already visible in two dimensions.
For the shear \(T_t(x,y)=(x+ty,y)\), conditioning on \(y\) and applying
Plancherel gives the exact formula
\begin{equation}
 h^2((T_t)_\#\mu_{\alpha,2},\mu_{\alpha,2})
 =2\int_\R(1-e^{-|t\xi|^\alpha})
                  |\widehat g_\alpha(\xi)|^2\,d\xi.
 \label{eq:stable-shear-hellinger-formula}
\end{equation}
Consequently
\[
 h^2((T_t)_\#\mu_{\alpha,2},\mu_{\alpha,2})
 \sim2|t|^\alpha\int_\R|\xi|^\alpha
                         |\widehat g_\alpha(\xi)|^2\,d\xi,
 \qquad t\to0,
\]
whereas a diagonal scale perturbation has squared Hellinger cost of order
\(t^2\).
Thus anisotropic stable Hellinger coercivity is not generated by a quadratic
form from one finite-dimensional matrix Fisher operator.

In particular, Hilbert--Schmidt proximity is necessary but no longer
sufficient.
On independent two-coordinate blocks, let
\[
 T_a(x_k,y_k)_{k\ge1}=(x_k+a_ky_k,y_k)_{k\ge1}.
\]
Then \(T_a-I\in\mathcal S_2\) exactly when \(a\in\ell^2\), while
\cref{thm:complete-criterion-symmetric-stable-coordinates} gives
\((T_a)_\#\mu_\alpha\sim\mu_\alpha\) exactly when
\(a\in\ell^\alpha\).
Any \(a\in\ell^2\setminus\ell^\alpha\) therefore gives a
Hilbert--Schmidt perturbation with singular image law, even though the
stable coordinate density is everywhere positive and has finite location
and scale Fisher information.
Conversely, when one source is mixed into many outputs, the criterion
charges the \(\ell^2\)-norm of that whole column only once.
An entrywise \(\ell^\alpha\) condition would therefore be too strong.
The correct stable criterion is the anisotropic mixing cost
\eqref{eq:stable-mixing-cost}, not a Schatten-class replacement for the
finite-variance theorem.
